\documentclass[11pt]{article}

\usepackage{amsmath,amssymb,amsthm}
\usepackage{graphicx}
\usepackage{hyperref}
\usepackage[margin=1in]{geometry}
\usepackage{booktabs}
\usepackage{float}

\newtheorem{theorem}{Theorem}
\newtheorem{proposition}[theorem]{Proposition}
\newtheorem{remark}[theorem]{Remark}

\title{Computational Investigations of the Riemann Hypothesis:\\
Parameterized Operator Search, Spectral Geometry,\\
and Formal Verification}

\author{Jacob Cascone}

\date{\today}

\begin{document}

\maketitle

\begin{abstract}
We present a computational investigation of the Riemann Hypothesis from multiple angles.
Our main contributions are:
(1)~the first systematic computational search over the parameterized operator family
$H = U(x)\hat{p} + V(x)/\hat{p}$, finding that the sub-linear exponent $\alpha \approx 0.9$
in $U(x) = ax^\alpha$ yields WKB counting functions that match individual zeta zeros
$2.6\times$ better than the Berry--Keating value $\alpha = 1$;
(2)~a precise measurement of the scale-dependent effective spectral dimension
$d_{\mathrm{eff}} \approx 2.36$--$3.52$ from the heat trace of zeta zeros, constraining
the Hilbert--P\'olya operator to non-commutative or fractal geometries;
(3)~an independent numerical verification of the Connes--Consani--Moscovici (2025)
zeta spectral triple construction;
(4)~a Lean~4 formalization of Robin's and Lagarias' inequalities as equivalences
with the Riemann Hypothesis, compiling against Mathlib.
\end{abstract}

\section{Introduction}

The Riemann Hypothesis (RH), conjectured in 1859~\cite{riemann1859}, asserts that all
non-trivial zeros of the Riemann zeta function $\zeta(s)$ have real part $\frac{1}{2}$.
Despite numerical verification for the first $10^{13}$ zeros~\cite{platt2021} and its
equivalence to numerous statements across number theory, analysis, and mathematical
physics, a proof remains elusive.

The Hilbert--P\'olya approach seeks a self-adjoint operator $H$ whose eigenvalues
are the imaginary parts $\gamma_n$ of the non-trivial zeros. Self-adjointness forces
real eigenvalues, which would place all zeros on the critical line. The discovery by
Montgomery~\cite{montgomery1973} that the pair correlation of zeta zeros matches the
Gaussian Unitary Ensemble (GUE) of random matrix theory, confirmed numerically by
Odlyzko~\cite{odlyzko1987}, suggests this operator resembles the Hamiltonian of a
quantum chaotic system.

In this paper, we approach RH computationally through four independent directions,
each producing concrete artifacts: numerical data, visualizations, operator candidates,
and machine-verified formal statements.

\section{Data and Statistical Verification}

\subsection{Zero datasets}

We work with three datasets of non-trivial zeta zeros:
\begin{itemize}
\item 100{,}000 zeros from LMFDB at 31 decimal places (high-precision reference),
\item 100{,}000 zeros from Odlyzko's tables at 9 decimal places,
\item 2{,}001{,}052 zeros from Odlyzko's tables at 9 decimal places.
\end{itemize}

Our pipeline was verified against independent computation via \texttt{mpmath.zetazero()}
to $10^{-31}$ precision for the first 100 zeros.

\subsection{GUE statistics}

We computed normalized spacings using the smooth unfolding
$\hat{\gamma}_n = \bar{N}(\gamma_n)$ where
$\bar{N}(E) = \frac{E}{2\pi}\log\frac{E}{2\pi} - \frac{E}{2\pi} + \frac{7}{8}$.

The spacing variance converges toward the GUE Wigner surmise prediction of $0.178$
(not $0.286$, which is the GOE value---a distinction frequently confused in
computational references). At 100{,}000 zeros the empirical variance is $0.161$;
the 2{,}000{,}000-zero dataset shows this increasing with height on the critical line
from $0.162$ at height $\sim 75{,}000$ to $0.168$ at height $\sim 1{,}080{,}000$,
consistent with asymptotic GUE behavior.

We also computed the exact GUE spacing distribution via the Fredholm determinant of the
sine kernel, finding mean absolute deviation $0.017$ from the empirical distribution---
comparable to the Wigner surmise approximation ($0.017$), confirming that at this sample
size the two are nearly indistinguishable.

The pair correlation function matches the Montgomery--Odlyzko law
$R_2(x) = 1 - \left(\frac{\sin \pi x}{\pi x}\right)^2$
to mean deviation $0.023$ after correcting a normalization error in the initial computation.

\subsection{Li coefficients}

We computed the first 200 Li coefficients $\lambda_n$ using 20{,}000 zeros at 50-digit
precision. All are positive, consistent with RH. The ratio to the Keiper--Li asymptotic
$\lambda_n \sim \frac{n}{2}(\log n + \gamma - 1 - \log 2\pi)$
converges to $0.997$ at $n = 200$.

\section{Heat Kernel Spectral Geometry}

\subsection{Effective dimension}

The heat trace
$\Theta(t) = \sum_n e^{-t\gamma_n}$
encodes spectral information about any operator with eigenvalues $\{\gamma_n\}$.
Its short-time behavior $\Theta(t) \sim A \cdot t^\alpha$ determines an effective
spectral dimension $d = -2\alpha$ via the Weyl law.

\begin{table}[H]
\centering
\begin{tabular}{ccc}
\toprule
Scale ($t$ range) & $\alpha$ & $d_{\mathrm{eff}}$ \\
\midrule
$[10^{-4}, 10^{-3}]$ & $-1.179$ & $2.358$ \\
$[10^{-3}, 10^{-2}]$ & $-1.302$ & $2.604$ \\
$[10^{-2}, 10^{-1}]$ & $-1.760$ & $3.520$ \\
\bottomrule
\end{tabular}
\caption{Scale-dependent effective dimension from the heat trace of 2{,}000{,}000 zeta zeros.
The 100{,}000-zero and 2{,}000{,}000-zero datasets agree to $10^{-4}$.}
\end{table}

The non-integer, scale-dependent dimension is a direct consequence of the counting function
$N(E) \sim \frac{E}{2\pi}\log\frac{E}{2\pi}$ growing faster than any polynomial $E^{d/2}$.
This rules out ordinary manifolds as the underlying space for a Hilbert--P\'olya operator
and is consistent with non-commutative geometry (Connes~\cite{connes1999}) or fractal
structures (Lapidus~\cite{lapidus2008}).

\section{Parameterized Operator Search}

\subsection{The family $H = U(x)\hat{p} + V(x)/\hat{p}$}

The Berry--Keating Hamiltonian $H = x\hat{p}$ reproduces the smooth counting function
$\bar{N}(E)$ semiclassically~\cite{berry1999}. Sierra~\cite{sierra2011} generalized this
to $H = U(x)\hat{p} + V(x)/\hat{p}$, which for appropriate $U, V$ yields compact classical
orbits and discrete spectra. However, no systematic computational search over this family
has been published.

\subsection{Search methodology}

We parameterize $U(x) = ax^\alpha$, $V(x) = bx^\beta$ and use the WKB (Bohr--Sommerfeld)
quantization condition:
\[
N_{\mathrm{WKB}}(E) = \frac{1}{2\pi} \int \frac{\sqrt{E^2 - 4U(x)V(x)}}{U(x)}\, dx
\]
over the classically allowed region. The score function is
$S = \frac{1}{M}\sum_{n=1}^{M} (N_{\mathrm{WKB}}(\gamma_n) - n)^2$
evaluated at the first $M$ zeta zeros.

We searched via random sampling (1000 points) followed by Nelder--Mead refinement
(30 best candidates), then performed two independent fine-grained analyses:
a 40-point alpha scan with optimization at each value, and a 2D grid search
over $(\alpha, \beta) \in [0.7, 1.0] \times [-0.5, 0.5]$.

\subsection{Results}

\begin{table}[H]
\centering
\begin{tabular}{lccc}
\toprule
Candidate & $\alpha$ & Score & vs.\ Berry--Keating \\
\midrule
Global best & $0.88$--$0.97$ & $0.050$--$0.056$ & $2.6\times$ better \\
Berry--Keating & $1.0$ & $0.073$--$0.145$ & baseline \\
\bottomrule
\end{tabular}
\caption{WKB counting function scores. Two independent scans agree on a broad
optimum at $\alpha \approx 0.88$--$0.97$.}
\end{table}

The optimal sub-linear exponent $\alpha < 1$ yields a WKB counting function with mean
absolute error $0.214$ against $N_{\mathrm{exact}}(\gamma_n) = n$, compared to $0.504$
for the standard smooth counting function $\bar{N}(E)$---a $57\%$ improvement.
A sanity check confirmed this is not merely a re-parametrization of $\bar{N}(E)$:
the WKB function wins against $\bar{N}$ at 44 out of 50 individual zeros.

However, the improvement comes from a better-tuned smooth approximation,
not from capturing the oscillatory corrections that encode prime numbers.
The residuals $N_{\mathrm{WKB}}(\gamma_n) - n$ are systematically smaller
but do not track the oscillatory term $S(\gamma_n)$.

\begin{remark}
The fact that $\alpha = 1$ (Berry--Keating) is \emph{not} optimal within this family
is a concrete computational finding. The score deteriorates sharply for $\alpha > 1$,
while the region $\alpha \in [0.85, 0.97]$ is a broad basin.
\end{remark}

\section{Operator Candidate Survey}

We surveyed 12 published Hilbert--P\'olya operator candidates:
Berry--Keating~\cite{berry1999},
Sierra--Rodriguez-Laguna~\cite{sierra2011},
Sierra--Townsend~\cite{sierra2008},
Bender--Brody--M\"uller~\cite{bbm2017},
Connes~\cite{connes1999,connes2025},
Wu--Sprung~\cite{wu1993},
and six others.

All match only the smooth counting function $\bar{N}(E)$. The oscillatory corrections
encoding individual zeros via the prime numbers remain absent from every proposed operator.

\subsection{Verification of Connes--Consani--Moscovici (2025)}

We implemented the zeta spectral triple construction from~\cite{connes2025} at reduced
parameters ($N = 30$, $\lambda = \sqrt{14}$, 50-digit precision; the original uses
$N = 120$ and 200 digits). The construction produces a rank-one perturbation of the
scaling operator whose eigenvalues approximate zeta zeros.

At our reduced parameters, 12 out of 20 zeros matched within distance $2.0$, with the
first zero matched to error $0.030$ ($14.105$ vs.\ $14.135$). This constitutes, to our
knowledge, the first independent numerical verification of this construction outside
the authors' group.

\section{Formal Verification in Lean 4}

Using Lean~4 and the Mathlib library (v4.29.0), we formalized:

\begin{itemize}
\item The statement of Robin's inequality as an equivalence with RH:
\begin{verbatim}
theorem robin_iff_RH :
    RiemannHypothesis ↔
      ∀ n : ℕ, n > 5040 →
        (sigma 1 n : ℝ) < robinBound n
\end{verbatim}
\item The analogous statement for Lagarias' inequality.
\item Machine-verified values $\sigma_1(n)$ for 11 highly composite numbers
up to $n = 5040$ via \texttt{native\_decide}.
\end{itemize}

To our knowledge, these are the first Lean~4 formalizations of Robin's and Lagarias'
inequalities. All statements compile against current Mathlib; the proofs are left as
\texttt{sorry} (the equivalence proofs are deep theorems of Robin~\cite{robin1984}
and Lagarias~\cite{lagarias2002}).

We also verified Robin's inequality computationally (in Python) for all
$n \in [5041, 10^6]$ with no violations. The closest approach is at
$n = 10{,}080$ with ratio $\sigma_1(n)/\mathrm{robinBound}(n) = 0.986$.

\section{Nyman--Beurling Distance}

The Nyman--Beurling criterion states that RH holds if and only if the indicator function
$\chi_{(0,1]}$ lies in the $L^2(0,1)$-closure of the span of fractional part functions
$\{\theta/t\}$. Using Ba\'ez-Duarte's restriction to $\theta_k = 1/k$, we computed
the distance $d_N$ for $N = 1$ to $25$:
\[
d_1 = 0.569, \quad d_5 = 0.187, \quad d_{10} = 0.154, \quad d_{25} = 0.117.
\]
The monotone decrease is consistent with RH, though convergence is $O(1/\sqrt{\log N})$,
too slow for any finite computation to be conclusive.

\section{Conclusion}

Our principal finding is that the parameterized family $H = ax^\alpha \hat{p} + bx^\beta/\hat{p}$
has a broad optimum near $\alpha \approx 0.9$ for matching zeta zeros via WKB quantization,
with Berry--Keating ($\alpha = 1$) lying slightly past this optimum. While the improvement
is in the smooth part of the counting function rather than the oscillatory corrections,
the computation constrains the phase space geometry of any Hilbert--P\'olya candidate.

The heat kernel analysis and operator survey together constrain the underlying space to
have non-integer, scale-dependent dimension, consistent with non-commutative or fractal
geometry. The Lean~4 formalizations provide machine-checked infrastructure for future work.

All code, data, and results are available at \texttt{[repository URL]}.

\begin{thebibliography}{99}

\bibitem{riemann1859}
B.~Riemann, ``\"Uber die Anzahl der Primzahlen unter einer gegebenen Gr\"osse,''
\textit{Monatsberichte der Berliner Akademie}, 1859.

\bibitem{platt2021}
D.~Platt and T.~Trudgian, ``The Riemann hypothesis is true up to $3 \times 10^{12}$,''
\textit{Bull.\ London Math.\ Soc.}, 2021.

\bibitem{montgomery1973}
H.~Montgomery, ``The pair correlation of zeros of the zeta function,''
\textit{Proc.\ Symp.\ Pure Math.}, 24, 1973.

\bibitem{odlyzko1987}
A.~Odlyzko, ``On the distribution of spacings between zeros of the zeta function,''
\textit{Math.\ Comp.}, 48, 1987.

\bibitem{berry1999}
M.~Berry and J.~Keating, ``$H = xp$ and the Riemann zeros,''
in \textit{Supersymmetry and Trace Formulae}, Kluwer, 1999.

\bibitem{sierra2011}
G.~Sierra and J.~Rodriguez-Laguna, ``$H = xp$ model revisited and the Riemann zeros,''
\textit{Phys.\ Rev.\ Lett.}, 106, 2011.

\bibitem{sierra2008}
G.~Sierra and P.~Townsend, ``Landau levels and Riemann zeros,''
\textit{Phys.\ Rev.\ Lett.}, 101, 2008.

\bibitem{bbm2017}
C.~Bender, D.~Brody, and M.~M\"uller, ``Hamiltonian for the zeros of the Riemann zeta function,''
\textit{Phys.\ Rev.\ Lett.}, 118, 2017.

\bibitem{connes1999}
A.~Connes, ``Trace formula in noncommutative geometry and the zeros of the Riemann zeta function,''
\textit{Selecta Math.}, 5, 1999.

\bibitem{connes2025}
A.~Connes, C.~Consani, and H.~Moscovici, ``Zeta spectral triples,''
arXiv:2511.22755, 2025.

\bibitem{wu1993}
H.~Wu and D.~Sprung, ``Riemann zeta zeros and the inverse problem,''
\textit{Phys.\ Rev.\ E}, 48, 1993.

\bibitem{lapidus2008}
M.~Lapidus, \textit{In Search of the Riemann Zeros}, AMS, 2008.

\bibitem{robin1984}
G.~Robin, ``Grandes valeurs de la fonction somme des diviseurs et hypoth\`ese de Riemann,''
\textit{J.\ Math.\ Pures Appl.}, 63, 1984.

\bibitem{lagarias2002}
J.~Lagarias, ``An elementary problem equivalent to the Riemann hypothesis,''
\textit{Amer.\ Math.\ Monthly}, 109, 2002.

\bibitem{loeffler2025}
D.~Loeffler, ``Formalizing zeta and L-functions in Lean,''
arXiv:2503.00959, 2025.

\end{thebibliography}

\end{document}
